\section{Triton softmax：row fits in one block}

Source coverage: \texttt{triton\_softmax\_example}. 本节覆盖的核心概念包括：row-wise reduction, max subtraction, exp, normalization。


\begin{figure}[H]
\centering
\includegraphics[width=0.86\textwidth]{triton-softmax.png}
\caption{triton softmax.png}
\figsource{CS336 Spring 2026 Lecture 6 官方可执行讲义。}
\end{figure}
\begin{knowledgebox}{读图：triton softmax.png}
这张图用于把抽象的 kernel 代码连接到硬件执行。读图时先看数据位于 HBM、shared memory 还是 registers，再看线程/blocks 如何映射到 SM。性能优化的关键是让数据在更靠近计算单元的位置被复用，减少慢速 HBM 往返。
\end{knowledgebox}

\[
\mathrm{softmax}(x_i)=\frac{e^{x_i-m}}{\sum_j e^{x_j-m}},\qquad m=\max_j x_j.
\]
\begin{importantbox}{fused softmax 的资源意义}
Naive softmax 需要 max、subtract、exp、sum、divide 多个 pass。若一整行能放进一个 Triton block，就可以在 shared memory/register 中完成 reduction，只读写 HBM 各一次左右。
\end{importantbox}
\subsection{本章小结}
本节说明了一个 kernel optimization 的共同模式：先确认语义正确，再测时间，再看 profiler，最后用 block、shared memory、tiling、fusion 或 layout 调整降低数据移动。

